\section{Solution}

We answer Question~\ref{1.2} in the affirmative by constructing an explicit
computably represented structure with the required properties.

\begin{theorem}\label{main}
There is a countable computably-represented structure $\A$ such that:
\begin{enumerate}
\item[(a)] $\A$ is computably $\AUT$-countable on a cone; and
\item[(b)] for every finite tuple $\bar a$ of parameters from $\A$ there is an
automorphism $\pi$ of $\A$ such that $\pi$ is not $\Sigma^{in}_1$-definable in
$\A$ relative to the parameters $\bar a\cup\pi(\bar a)$.
\end{enumerate}
\end{theorem}

\subsection{The structure}

We work in a language with a single binary relation symbol $E$ (a graph) together
with finitely many unary ``colour'' predicates that we describe along the way; all
of these can be coded into a single binary relation, so the structure is a graph,
and our explicit description gives a computable representation.

\medskip

\noindent\textbf{Building blocks.}
For each $n\in\N$ we build a finite connected ``gadget'' $G_n$ whose automorphism
group is the cyclic group $\Z/2\Z$, realized by a single nontrivial automorphism
$\sigma_n$ of order $2$. We do this so that:
\begin{enumerate}
\item[(i)] $G_n$ is rigid except for the involution $\sigma_n$;
\item[(ii)] $\sigma_n$ swaps two distinguished vertices $p_n,q_n$ of $G_n$
(its ``poles''), $\sigma_n(p_n)=q_n$;
\item[(iii)] every vertex $v$ of $G_n$ is moved by $\sigma_n$, i.e.\ $\sigma_n$ is
fixed-point-free; and
\item[(iv)] $G_n$ uniformly contains an internal ``distance gadget'' so that the
following holds: in any representation $\B\cong\A$, in order to compute, on input a
vertex $x\in G_n$, the value $\sigma_n(x)$, one must decide membership in a fixed
$\Sigma^{in}_1$-but-not-uniformly-computable relation. We make this precise below
by encoding a "long path" obstruction. (See the construction of $H_k$ below.)
\end{enumerate}

A concrete realization of (i)--(iii): let $G_n$ be the graph on vertex set
$\{p_n,q_n\}\cup\{a_{n,j},b_{n,j}: 0\le j<3\}$ obtained as follows. Put a path
$p_n - a_{n,0} - a_{n,1} - a_{n,2}$ of length $3$, a path
$q_n - b_{n,0} - b_{n,1} - b_{n,2}$ of length $3$, and add the single edge
$a_{n,2}-b_{n,2}$ joining the two endpoints. The map $\sigma_n$ with
$\sigma_n(p_n)=q_n$, $\sigma_n(a_{n,j})=b_{n,j}$ (and its inverse) is the unique
nontrivial automorphism: the graph is a path of $8$ vertices, whose only nontrivial
automorphism is the order-$2$ reflection, which is fixed-point-free and swaps the
two endpoints $p_n,q_n$. Thus $\AUT(G_n)=\{\mathrm{id},\sigma_n\}\cong\Z/2\Z$.

\medskip

\noindent\textbf{Encoding noncomputability into the action.}
Fix a set $K\subseteq\N$ that is computably enumerable but not computable (e.g.\ the
halting set). We attach to each gadget $G_n$ a uniformly given ``marker''. Choose
two non-isomorphic finite rigid graphs $M^0$ and $M^1$ (markers), say $M^0$ a
triangle and $M^1$ a $4$-cycle; these are rigid up to the constraint that we only
allow them as labelled attachments (we attach $M^0$ or $M^1$ to a vertex by an edge
to a distinguished apex of the marker, and we use a colour predicate to mark the
attachment vertex of the marker so that markers admit no nontrivial automorphism
in context). For each $n$, attach a copy of $M^{e_n}$ to the pole $p_n$ and a copy of
$M^{1-e_n}$ to the pole $q_n$, where $e_n=1$ if $n\in K$ and $e_n=0$ otherwise.

This attachment is asymmetric: $p_n$ carries marker $M^{e_n}$ and $q_n$ carries
$M^{1-e_n}$. Crucially the involution $\sigma_n$ continues to be an automorphism of
$G_n$ with markers \emph{only if} the markers attached at $p_n$ and $q_n$ are
swapped consistently; since $\sigma_n(p_n)=q_n$, the marker at $p_n$ is sent to the
position of $q_n$, so $\sigma_n$ is an automorphism of the marked gadget iff the
marker types are swapped under $\sigma_n$. Because $M^{e_n}$ and $M^{1-e_n}$ are
different graphs, $\sigma_n$ is \emph{not} an automorphism of the marked gadget.
This would destroy the involution, which we do not want; we therefore use markers
differently, as follows.

\medskip

\noindent\textbf{Correct encoding: hide the orientation, not destroy the symmetry.}
We instead attach to \emph{both} poles a single common marker which records,
\emph{abstractly}, an unordered pair, but place an additional ``orientation tag''
that is computably accessible only via a $\Sigma^{in}_1$ (but not computable)
search. Concretely:

\begin{itemize}
\item Keep $G_n$ (the $8$-vertex path) with its involution $\sigma_n$
swapping $p_n\leftrightarrow q_n$.
\item Add to $G_n$ an additional vertex $c_n$ adjacent to \emph{both} $p_n$ and
$q_n$ (a common ``hub''). Now $\sigma_n$ extends to fix $c_n$ and remains the
unique nontrivial automorphism (the symmetry of the cycle thus formed is still the
single reflection).
\item Attach to $c_n$ a long disjoint pendant path $P_n$ of length $\ell_n+10$,
where $\ell_n$ is the (finite, since $K$ is c.e.) stage at which $n$ enters $K$ if
$n\in K$, and where, if $n\notin K$, we instead attach to $c_n$ an \emph{infinite}
ray $R_n$ (a one-way infinite path).
\end{itemize}

The pendant attached at $c_n$ is fixed setwise by every automorphism (it is the
unique pendant path/ray at $c_n$), and it does not affect the existence of
$\sigma_n$. Thus in the gadget $G_n^\ast$ (gadget plus hub plus pendant) we still
have $\AUT(G_n^\ast)=\{\mathrm{id},\sigma_n\}$, with $\sigma_n$ swapping
$p_n\leftrightarrow q_n$ and fixing $c_n$ and the pendant pointwise. The data
``$n\in K$ vs.\ $n\notin K$'' is encoded in whether the pendant is finite or
infinite, and \emph{this difference is detectable by a $\Sigma^{in}_1$ search}
(``does the pendant terminate?'') but not computably.

\medskip

\noindent\textbf{The whole structure.}
Let $\A$ be the disjoint union $\bigsqcup_{n\in\N} G_n^\ast$, together with a unary
predicate $U$ marking each hub vertex $c_n$ and a unary predicate $P$ marking each
pole vertex $p_n,q_n$. The disjoint pieces are distinguished by connectivity. This
$\A$ is computably represented: we list vertices, edges, and the predicates
uniformly in $n$ and $j$; the only ``infinite'' object is the ray $R_n$ for
$n\notin K$, but its vertices and edges are enumerated computably (we always lay
out a ray, and at the finite stage $\ell_n$ at which $n$ would enter $K$ — for those
$n$ that do enter — we never actually do, since we decided in advance; precisely, we
fix once and for all the standard enumeration of $K$ and define the pendant as the
ray with a ``cap'' vertex added at position $\ell_n$ if $n$ enters $K$ at stage
$\ell_n$). All of this is uniformly computable. Hence $\A$ has a computable
representation.

\subsection{Automorphism group and proof of (a)}

\begin{lemma}\label{autgroup}
For every representation $\B\cong\A$, $\AUT(\B)\cong\prod_{n\in\N}\Z/2\Z$, where the
$n$-th factor is generated by the involution $\tau_n$ that acts as $\sigma_n$ on the
copy of $G_n^\ast$ and as the identity elsewhere.
\end{lemma}

\begin{proof}
Work in $\A$ itself; the statement is isomorphism-invariant. Each connected
component of $\A$ is some $G_n^\ast$. An automorphism of $\A$ permutes the
connected components; but the components are pairwise non-isomorphic: $G_n^\ast$ and
$G_m^\ast$ for $n\neq m$ differ either by the pendant length (when both are in $K$
or one is) or by finiteness vs.\ infiniteness of the pendant. Indeed we may arrange
all pendant lengths $\ell_n+10$ to be distinct (and all $\geq 11$), and the
no-$K$ ray case to be the unique infinite case for each component, so that
distinct components are non-isomorphic. Hence every automorphism fixes each
component setwise, and restricts to an automorphism of that component. Since
$\AUT(G_n^\ast)=\{\mathrm{id},\sigma_n\}$, the global automorphism group is the
restricted product\,/\,product $\prod_n\{\mathrm{id},\sigma_n\}\cong(\Z/2\Z)^\N$, with
each coordinate chosen independently.
\end{proof}

\begin{lemma}[Proof of (a)]\label{conemany}
$\A$ is computably $\AUT$-countable on a cone.
\end{lemma}

\begin{proof}
Let $C=K$ (the halting set), and let $\B\cong\A$ be arbitrary. We must show every
automorphism $\pi$ of $\B$ is computable from $\B\oplus K$.

By Lemma~\ref{autgroup}, $\pi$ fixes each component setwise; on the component
isomorphic to $G_n^\ast$, $\pi$ acts either as the identity or as the involution
$\sigma_n$. Given a vertex $v\in\B$, we compute $\pi(v)$ as follows. From $\B$ we
compute the connected component $C_v$ of $v$; this component is a copy of some
$G_n^\ast$. Using $\B$ we identify the local structure: the predicates $U,P$ and the
graph metric let us, computably in $\B$, locate the hub $c$, the two poles
$p,q$, the path skeleton, and the pendant attached at the hub.

The only nontrivial thing $\pi$ can do on $C_v$ is the involution. Whether
$\pi\restriction C_v$ is the identity or the involution is a single bit
$b(C_v)\in\{0,1\}$; if we know this bit, then $\pi(v)$ is computable from $\B$
(apply $\sigma$ or $\mathrm{id}$ according to the local labelling). So it suffices to
show $b(\cdot)$ is computable from $\B\oplus K$. But $b(C_v)$ is determined by
$\pi$, and we are merely \emph{computing} the given $\pi$, so we may consult $\pi$:
to determine $b(C_v)$, compute $\pi(p)$ where $p$ is the (computably located) pole
of $C_v$; if $\pi(p)=p$ then $b=0$, else $b=1$. Computing $\pi(p)$ requires querying
$\pi$, which is allowed since $\pi$ is the function we are computing. Indeed the
function $v\mapsto\pi(v)$ is what we must compute; the reduction is: locate the pole
$p$ of $C_v$ (computable in $\B$), the value $\pi(p)$ is one of the two poles
$\{p,q\}$ and is part of the graph of $\pi$. The point is to reduce computing
$\pi$ everywhere to computing it on the two poles of each component, plus the
local computable structure. We make this honest by the following: $\pi$, as a set
of pairs, restricted to the poles, is a function
$\rho:\{p_n,q_n:n\}\to\{p_n,q_n:n\}$ with $\rho(p_n),\rho(q_n)\in\{p_n,q_n\}$. We
claim $\rho\le_T\B\oplus K$ — but in fact $\rho$ \emph{is} $\le_T \pi$ trivially and
we want $\pi\le_T\B\oplus K$.
\end{proof}

The previous paragraph contains the substantive point and we now state it cleanly,
since $\pi$ ranges over \emph{all} automorphisms and ``computability of $\pi$''
means there is a single oracle Turing functional computing $\pi$ from $\B\oplus K$.

\begin{proof}[Proof of Lemma~\ref{conemany}, completed]
Let $\B\cong\A$ and $\pi\in\AUT(\B)$. We give a $\B\oplus K$-effective procedure that
computes the graph of $\pi$. Recall $\AUT(\B)=(\Z/2\Z)^\N$, so $\pi$ is determined by
the bit-sequence $(b_n)_n$, $b_n\in\{0,1\}$, where $b_n=1$ iff $\pi$ acts as the
involution on the $n$-th component.

The graph of $\pi$ as a \emph{set} is automatically computable from $\B\oplus K$
\emph{provided} the map $C\mapsto b(C)$ (component $\mapsto$ its bit) is computable
from $\B\oplus K$ — but this map is not determined by $\B$; it is determined by
$\pi$. The correct interpretation of Definition~\ref{1.1} is that we are given the
oracle $\B$ and we must compute the \emph{specific} automorphism $\pi$; thus the
data of $\pi$ is presented to us through the requirement, and ``$\pi$ is computable
from $\B\oplus K$'' means: the function $v\mapsto\pi(v)$ is Turing-reducible to
$\B\oplus K$. Equivalently $\pi\le_T \B\oplus K$ as a function/oracle.

Now fix the particular $\pi$. Its bit-sequence $(b_n)$ is an arbitrary element of
$2^\N$, hence \emph{a priori} not computable from anything. We must therefore show
that $(b_n)\le_T \B\oplus K$ for \emph{this} $\pi$ — which is false in general since
$(b_n)$ can be any real! So we have shown $\A$ is \emph{not} computably
$\AUT$-countable on a cone with this construction.
\end{proof}

The honest analysis above shows that the disjoint-union construction \emph{fails}
condition (a): with infinitely many independent involutions the automorphism group
is uncountable, contradicting the standing hypothesis (Theorem~\ref{2.2}) that
$\AUT(\A)$ be countable. We must therefore redesign so that $\AUT(\A)$ is countable
while retaining the parameter-defeating behaviour.

\subsection{Corrected construction with countable automorphism group}

We now give a construction with $\AUT(\A)$ countable (indeed cyclic of order $2$ on
each "block" but globally controlled) that satisfies both (a) and (b).

\medskip

\noindent\textbf{Idea.} Use a single global involution that simultaneously swaps two
``halves'' of the structure, where to recover the swap on far-away parts of the
structure one needs unboundedly much $K$-information, so that no finite parameter
tuple makes the involution $\Sigma^{in}_1$-definable, yet on a cone the whole
involution is computable because $K$ is available as a uniform oracle.

\medskip

\noindent\textbf{Construction.} Let $\A$ consist of two copies of $\N$, called
$L=\{\ell_n:n\in\N\}$ ("left") and $R=\{r_n:n\in\N\}$ ("right"), with the following.
For each $n$, attach to $\ell_n$ and to $r_n$ identical finite "address gadgets"
$A_n$ that uniquely encode the index $n$ rigidly (e.g.\ a rigidly-coloured path of
length $n+3$ hung off the vertex, using two colours to write $n$ in binary along
the path; these are rigid). This makes $\ell_n$ and $r_n$ each definable without
parameters by their address. Add a unary predicate $L$-colour on all of $L$ and an
$R$-colour on all of $R$. Finally, for each $n$, encode the bit $e_n=[\,n\in K\,]$
by attaching to $\ell_n$ a pendant triangle if $e_n=1$ and a pendant square if
$e_n=0$, and to $r_n$ the \emph{opposite} marker (square if $e_n=1$, triangle if
$e_n=0$).

The structure $\A$ has exactly two automorphisms: the identity, and the global
involution $\pi$ defined by $\pi(\ell_n)=r_n$, $\pi(r_n)=\ell_n$ for all $n$
(extended to gadgets in the unique way). Indeed: any automorphism preserves the
address gadget at each vertex, hence maps the unique vertex of address $n$ to a
vertex of address $n$; the only two such vertices are $\ell_n,r_n$. The $L/R$
colours are \emph{not} part of the language (we omit them, otherwise the involution
would fail): we omit $L$-/$R$-colours. Then an automorphism either fixes the pair
$\{\ell_n,r_n\}$ pointwise for all $n$ (identity) or swaps it for all $n$; but the
markers force consistency: at index $n$, the marker at $\ell_n$ (resp.\ $r_n$) is a
triangle/square depending on $e_n$, and the swap $\pi$ sends the marker at $\ell_n$
to the position of $r_n$, where the marker is the opposite shape — so $\pi$ would
need to turn a triangle into a square, which is impossible.

\medskip

This again kills the involution. The fundamental tension is now clear and is the
content of Theorem~\ref{2.3}: \emph{any} automorphism of a structure that is
computably $\AUT$-countable on a cone is, by Theorem~\ref{2.3}, $\Sigma^{in}_1$
definable from \emph{some} finite parameter tuple. Question~\ref{1.2} asks whether
the parameters can be made to depend on $\pi$ in an essential way, i.e.\ whether one
\emph{fixed} tuple can fail.

\begin{remark}
After the analysis above I am unable to produce a correct construction within the
constraints, and the two natural approaches both collapse: making the involutions
independent forces an uncountable automorphism group (violating the standing
countability hypothesis and condition (a)), while making a single global involution
forces, via the address gadgets, that the involution swaps each pair
$\{\ell_n,r_n\}$, and Theorem~\ref{2.3} then guarantees a finite parameter set
defining it. I therefore cannot certify a complete proof of Theorem~\ref{main}.
\end{remark}

\medskip

\noindent\emph{Honest conclusion.} The arguments above do not constitute a valid
proof of Theorem~\ref{main}, and I withdraw the claimed theorem.